\documentclass{article}
\usepackage{iclr2027_conference,times}
\usepackage[T1]{fontenc}
\usepackage{lmodern}
\usepackage{hyperref}
\usepackage{url}
\usepackage{booktabs}
\usepackage{microtype}
\usepackage{enumitem}
\usepackage{xcolor}
\usepackage{amsmath,amssymb}
\usepackage{amsthm}
\usepackage{mathtools}

\newtheorem{theorem}{Theorem}
\newtheorem{proposition}{Proposition}
\newtheorem{lemma}{Lemma}
\newtheorem{corollary}{Corollary}
\newtheorem{remark}{Remark}
\newcommand{\E}{\mathbb E}
\newcommand{\R}{\mathbb R}
\newcommand{\1}{\mathbf 1}
\newcommand{\Jac}{\operatorname{Jac}}
\newcommand{\CCdim}{\operatorname{CCdim}}
\newcommand{\affdim}{\operatorname{affdim}}
\newcommand{\supp}{\operatorname{supp}}

\title{Jaccard Needs Exponentially Many Convex Scores}

\author{\textbf{Pahan Dewasurendra} \\ Johns Hopkins University}

\hypersetup{pdftitle={Jaccard Needs Exponentially Many Convex Scores},pdfauthor={Pahan Dewasurendra},pdfkeywords={convex calibration dimension, Jaccard index, intersection over union, surrogate losses, multi-label prediction, lower bounds},pdfsubject={preprint}}
\begin{document}
\maketitle

\begin{abstract}
Jaccard similarity, also called intersection over union, and the F1 score are pointwise monotone transforms of one another. Their convex surrogate complexities are nevertheless exponentially different. For instance-wise multi-label prediction with $s$ labels, we prove \[ 2^{s-1}-1 \;\leq\; \CCdim(L^{\rm Jac}) \;\leq\; 2^s-1. \] Thus every convex surrogate calibrated for Jaccard on all conditional label distributions needs exponentially many real-valued scores. In contrast, a recent tight result gives convex calibration dimension $\Theta(s^2)$ for F1. The lower bound comes from one explicit conditional distribution. It is supported on all sets containing a fixed core label, weights a set with $d$ optional labels by $1/d!$, and makes exactly those $2^{s-1}$ reports Bayes optimal. Strict positive definiteness of the corresponding Jaccard block then collapses its feasible Bayes subspace to a point.

We also identify an algebraic phase transition. In the normalized-overlap family interpolating from Dice/F1 to Jaccard, the exact score-matrix rank is $s^2-s+2$ at the F1 endpoint but $2^s$ after any positive Jaccard correction. Finally, we reconcile our theorem with a NeurIPS 2024 claim that an $s$-score convex logistic surrogate is Bayes-consistent for every multi-label loss. A two-label rational example gives zero surrogate regret and Jaccard regret $1/10$. The proof in that work transfers probability between two arbitrary label vectors, but its additive softmax probabilities form only an $s$-dimensional product family and cannot realize that transfer. Our results separate practical restricted-distribution methods from distribution-free convex calibration and show that, for Jaccard, this distinction is unavoidable.
\end{abstract}

\section{Introduction}

An image mask, tag vector, or structured binary prediction can be represented as a subset of $s$ labels. Two standard scores for a report $B$ and outcome $A$ are \[ F_1(A,B)=\frac{2|A\cap B|}{|A|+|B|}, \qquad \Jac(A,B)=\frac{|A\cap B|}{|A\cup B|}. \tag{1}\label{eq:scores} \] Both equal one on identical nonempty sets, and we adopt the usual convention that both equal one on two empty sets. Pointwise they contain the same ordering information: \[ \Jac(A,B)=\frac{F_1(A,B)}{2-F_1(A,B)}. \tag{2}\label{eq:monotone} \] This close relationship makes the contrast in this paper surprising.

For a discrete loss, convex calibration dimension is the smallest $d$ for which some loss convex in a score vector $u\in\R^d$, together with some link from scores to reports, is calibrated on every conditional outcome distribution \citep{ramaswamy2016ccdim}. It asks how many real-valued outputs a distribution-free convex surrogate fundamentally needs. For F1, calibrated surrogates with $O(s^2)$ scores have been known since \citet{zhang2020fmeasure}. A new exact analysis proves the matching lower bound $\CCdim(L^{F_1})=\Theta(s^2)$ and identifies the exact score rank $s^2-s+2$ \citep{zhang2026exact}. We prove that Jaccard instead needs $\Theta(2^s)$ scores.

This question matters because IoU is used precisely when coordinatewise metrics such as Hamming loss are inadequate. Bayes-optimal IoU decisions are already combinatorial even when conditional label probabilities factorize \citep{nowozin2014iou}. Practical losses include the Lov\'asz hinge and its extensions \citep{yu2018lovasz,berman2018lovasz}. Recent work shows that the usual Lov\'asz-hinge link is inconsistent for Jaccard \citep{finocchiaro2025lovasz}. That result concerns one surrogate and link. Our lower bound concerns every convex surrogate and every link below an exponential dimension.

There is an apparent contradiction. \citet{mao2024multilabel} claim that a new $s$-score convex multi-label logistic loss is Bayes-consistent, with an excess-risk bound, for any multi-label loss. Jaccard is an explicit example in their theorem. Our lower bound rules this out. We resolve the contradiction constructively with an exact two-label counterexample and pinpoint the failed realizability step in their proof. The analogous full-output structured construction in \citet{mao2023structured} has one free score for each label vector, so its probability-transfer argument does not have this problem and is consistent with our exponential upper bound.

\paragraph{Contributions.}
\begin{enumerate}[leftmargin=*,itemsep=2pt,topsep=3pt]
\item We prove an exponential lower bound
$\CCdim(L^{\rm Jac})\geq 2^{s-1}-1$. The standard affine upper bound is $2^s-1$, determining the order exactly.
\item We give an explicit factorially weighted Bayes witness and a short
combinatorial identity showing that precisely all reports containing one core label tie. This produces exponentially many independent active Bayes constraints at one distribution.
\item We prove a sharp rank discontinuity for a continuous family of
normalized-overlap scores. The F1 endpoint has quadratic rank, while every positive Jaccard correction has full rank $2^s$.
\item We give a rational two-label falsification certificate for the claimed
general consistency of the low-dimensional multi-label logistic loss. It has zero surrogate regret but positive target regret, and does not depend on the empty-set convention.
\end{enumerate}

The lower bound is about expected per-instance Jaccard loss, also called the decision-theoretic formulation. It does not concern a ratio formed after aggregating confusion counts across an entire population. We return to this scope distinction in Section~\ref{sec:discussion}.

\section{Related work and position}
\label{sec:related}

\paragraph{Convex calibration dimension.}
\citet{ramaswamy2016ccdim} introduced convex calibration dimension for finite losses, proved affine and rank upper bounds, and developed the feasible-subspace lower bound used here. The flat-based elicitation framework of \citet{finocchiaro2021lower} recovers that theorem for minimizable convex surrogates and gives stronger global tools. Neither work analyzes Jaccard or IoU. An important distinction in both is that matrix rank supplies a generic upper bound, not a lower bound for arbitrary convex surrogates. This is why our full-rank theorem and factorial Bayes witness are separate contributions.

\paragraph{F-measures.}
F1 Bayes decisions can be recovered from a quadratic collection of conditional statistics. \citet{zhang2020fmeasure} turn this structure into convex calibrated surrogates. Most recently, \citet{zhang2026exact} prove the exact rank and affine dimension and construct a feasible-subspace witness giving the matching $\Theta(s^2)$ calibration dimension. Our witness is inspired by the same general question but differs in both weighting and geometry. It creates $2^{s-1}$ independent active comparisons rather than a quadratic family.

\paragraph{IoU optimization and the Lov\'asz hinge.}
\citet{nowozin2014iou} study Bayes decisions for IoU under a probabilistic model and derive an approximation for aggregate semantic segmentation. The Lov\'asz hinge and Lov\'asz-Softmax use submodularity to obtain practical convex objectives \citep{yu2018lovasz,berman2018lovasz}. For symmetric submodular losses, \citet{finocchiaro2022lovasz} identify the structured abstain target that the Lov\'asz hinge actually embeds. The recent label-dependent extension proves that the usual sign link is inconsistent for Jaccard when $s\geq3$ \citep{finocchiaro2025lovasz}. It leaves open whether a different convex surrogate or link of comparable dimension could be calibrated. Theorem~\ref{thm:main} rules out all of them.

\paragraph{Low-dimensional structured surrogates.}
The full-output structured logistic loss of \citet{mao2023structured} has a free score for each of the $2^s$ label vectors. The multi-label specialization of \citet{mao2024multilabel} replaces those scores by $s$ additive label scores and claims the same guarantee for every target loss. Section \ref{sec:counterexample} shows exactly why that reduction fails for Jaccard. Partial calibration under low-noise conditions and aggregation of several low-dimensional problems are studied by \citet{nueve2024tradeoff}. Those results complement our full-simplex lower bound.

Table~\ref{tab:separation} summarizes the sharp separation. The F1 entries are from \citet{zhang2026exact}; the Jaccard entries are proved here.

\begin{table}[t]
\centering
\caption{F1 and Jaccard are pointwise monotone-equivalent but have different
convex surrogate geometry.}
\label{tab:separation}
\begin{tabular}{lcc}
\toprule
Quantity on $s$ labels & F1/Dice & Jaccard/IoU\\
\midrule
Score-matrix rank & $s^2-s+2$ & $2^s$\\
Loss-column affine dimension & $s^2-s+1$ & $2^s-1$\\
Convex calibration dimension & $\Theta(s^2)$ & $\Theta(2^s)$\\
Explicit bounds & $\Omega(s^2)$ to $s^2-s+1$
                & $2^{s-1}-1$ to $2^s-1$\\
\bottomrule
\end{tabular}
\end{table}

\section{Setup and calibration dimension}
\label{sec:setup}

Let $[s]=\{1,\ldots,s\}$ and let outcomes and reports both range over $\mathcal Y=2^{[s]}$, of cardinality $N=2^s$. Define the Jaccard score matrix $K\in\R^{N\times N}$ and loss matrix $L^{\rm Jac}$ by \[
 K_{A,B}=\begin{cases}
  1,&A=B=\varnothing,\\
  |A\cap B|/|A\cup B|,&\text{otherwise},
 \end{cases}
\qquad L^{\rm Jac}=\1\1^\top-K. \tag{3}\label{eq:jacmatrix} \] Rows index outcomes and columns index reports. For a conditional distribution $q\in\Delta_N$, the Bayes reports maximize $q^\top K_{\cdot,B}$.

A surrogate is a function $\psi:\mathcal C\to\R_+^N$, where $\mathcal C\subseteq\R^d$ is convex and each coordinate $\psi_y$ is convex. It is calibrated for $L$ if it admits a link from $\mathcal C$ to reports such that, at every $q\in\Delta_N$, approaching optimal surrogate risk cannot retain a suboptimal target report. The convex calibration dimension $\CCdim(L)$ is the smallest such $d$ \citep{ramaswamy2016ccdim}. This definition permits arbitrary convex losses and arbitrary links.

For report $B$, its trigger polytope is \[ Q_B=\{q\in\Delta_N:q^\top L_{\cdot,B} \leq q^\top L_{\cdot,C}\ \text{for every }C\}. \tag{4}\label{eq:trigger} \] Let $\mu_{Q_B}(q)$ be the dimension of the largest linear subspace of two-sided feasible directions of $Q_B$ at $q$. The feasible-subspace theorem of \citet{ramaswamy2016ccdim} states \[ \CCdim(L)\geq \|q\|_0-\mu_{Q_B}(q)-1, \qquad q\in Q_B. \tag{5}\label{eq:ramalower} \] The newer flat-based framework recovers this bound and can strengthen it in other problems \citep{finocchiaro2021lower}. Our witness makes $\mu_{Q_B}(q)=0$, so the original theorem is already sharp enough.

\section{A full-rank phase transition}
\label{sec:rank}

We first isolate an algebraic distinction between F1 and Jaccard. For $0\leq\lambda<2$, define the normalized-overlap score \[ S_\lambda(A,B) =\frac{(2-\lambda)|A\cap B|} {|A|+|B|-\lambda|A\cap B|} \tag{6}\label{eq:family} \] for nonempty pairs, with value zero if exactly one set is empty and value one if both are empty. Then $S_0=F_1$ and $S_1=\Jac$. If $D=F_1(A,B)$, then \[ S_\lambda(A,B)=\frac{(2-\lambda)D}{2-\lambda D}, \tag{7}\label{eq:transform} \] so every member is a pointwise strictly increasing transform of F1.

\begin{theorem}[Rank phase transition]
\label{thm:rank}
Under the empty-empty convention above, the score matrix of $S_0$ has rank $s^2-s+2$. For every $\lambda\in(0,2)$, the score matrix of $S_\lambda$ is strictly positive definite and has rank $2^s$.
\end{theorem}

The endpoint rank is the exact F1 result of \citet{zhang2026exact}. We prove the positive-$\lambda$ statement. Write $x=|A\cap B|$, $a=|A|$, and $b=|B|$. For nonempty sets, \[ S_\lambda(A,B) =(2-\lambda)\int_0^\infty x e^{-t(a+b-\lambda x)}\,dt. \tag{8}\label{eq:integral} \] Let $\chi_A\in\{0,1\}^s$ be the incidence vector. Expanding $x e^{\lambda tx}$ into tensor powers gives, for coefficients $c_A$ on the nonempty sets,
\begin{align}
 c^\top S_\lambda c
 &=(2-\lambda)\int_0^\infty
 \sum_{r\geq1}\frac{(\lambda t)^{r-1}}{(r-1)!}
 \left\|\sum_{A\ne\varnothing}
 c_Ae^{-t|A|}\chi_A^{\otimes r}\right\|_2^2dt.
 \tag{9}\label{eq:tensor}
\end{align}
This is nonnegative. If it vanished, every displayed tensor sum would vanish. For any nonempty $T\subseteq[s]$, inspect a coordinate of the $|T|$th tensor that lists every element of $T$. It gives \[ \sum_{A\supseteq T}c_Ae^{-t|A|}=0. \tag{10}\label{eq:zeta} \] At any fixed $t>0$, these are all coordinates of the Boolean upper-zeta transform of $c_Ae^{-t|A|}$. The transform is invertible, so every $c_A=0$. The nonempty block is strictly positive definite. The empty-set convention adds an isolated $1\times1$ block, proving the theorem. A different positive definiteness proof for the Jaccard endpoint was given by \citet{bouchard2013jaccard}. Equation~\eqref{eq:tensor} is included because it covers the full family and makes the discontinuity transparent.

Full score rank alone is not a lower bound on arbitrary convex calibration dimension. It does imply that the $N$ score columns, and hence the $N$ loss columns, are affinely independent. Therefore \[ \affdim(L^{\rm Jac})=N-1. \tag{11}\label{eq:affdim} \] The generic affine construction of \citet{ramaswamy2016ccdim} now gives $\CCdim(L^{\rm Jac})\leq N-1$. The next section proves that no polynomial improvement is possible.

\section{The exponential calibration barrier}
\label{sec:barrier}

Fix one core label, denoted $\star$, and write $O=[s]\setminus\{\star\}$, with $n=s-1$. Consider the $2^n$ outcomes containing the core: \[ \mathcal U=\{\{\star\}\cup D:D\subseteq O\}. \tag{12}\label{eq:U} \] Define a distribution $p$ supported on $\mathcal U$ by \[ p_{\{\star\}\cup D} =\frac{1}{Z_n|D|!}, \qquad Z_n=\sum_{d=0}^n\binom nd\frac1{d!}. \tag{13}\label{eq:witness} \]

\begin{lemma}[Factorial tie]
\label{lem:tie}
Under $p$, precisely the reports in $\mathcal U$ maximize expected Jaccard similarity.
\end{lemma}

Here is the main identity. For $C\subseteq O$, ignore the common normalizer $Z_n$ and write the expected score of report $\{\star\}\cup C$ as \[ W(C)=\sum_{D\subseteq O}\frac1{|D|!} \frac{1+|C\cap D|}{1+|C\cup D|}. \tag{14}\label{eq:W} \] Then \[ W(C)=\sum_{d=0}^n\binom nd\frac1{(d+1)!}, \tag{15}\label{eq:constant} \] independent of $C$. To see why, add one optional label $j\notin C$ and pair $D=E$ with $D=E\cup\{j\}$. Put $a=|E\cap C|$, $b=|E\setminus C|$, and $c=|C|$. The paired increment is \[ \frac1{c+b+2} \left[ \frac1{(a+b+1)!} -\frac{a+1}{(a+b)!(c+b+1)} \right]. \tag{16}\label{eq:increment} \] For every fixed $b$, summing over the $a$ choices cancels by \[ \sum_{a=0}^c\binom ca\frac1{(a+b+1)!} =\frac1{c+b+1}\sum_{a=0}^c \binom ca\frac{a+1}{(a+b)!}. \tag{17}\label{eq:identity} \] Thus adding $j$ does not change $W$, and $W(C)=W(\varnothing)$ gives \eqref{eq:constant}. If a report omits $\star$, adding $\star$ strictly increases its intersection with every outcome in the support while leaving the union unchanged. Such a report is strictly suboptimal. Appendix \ref{app:tie} gives the complete cancellation proof.

\begin{theorem}[Exponential Jaccard calibration dimension]
\label{thm:main}
For every $s\geq2$, \[ 2^{s-1}-1\leq\CCdim(L^{\rm Jac})\leq2^s-1. \tag{18}\label{eq:main} \] In particular, $\CCdim(L^{\rm Jac})=\Theta(2^s)$.
\end{theorem}

\begin{proof}
The upper bound follows from \eqref{eq:affdim}. For the lower bound, let $r=|\mathcal U|=2^{s-1}$ and fix $B_0\in\mathcal U$. By Lemma~\ref{lem:tie}, $p\in Q_{B_0}$ and every comparison between $B_0$ and $B\in\mathcal U$ is active.

Restrict the score matrix to rows and columns in $\mathcal U$. It is a principal submatrix of the strictly positive definite Jaccard matrix, so its $r$ columns are linearly independent. Therefore the $r-1$ differences \[ K_{\mathcal U,B}-K_{\mathcal U,B_0}, \qquad B\in\mathcal U\setminus\{B_0\}, \tag{19}\label{eq:diffs} \] are linearly independent. Their common orthogonal complement in $\R^r$ is one-dimensional. The vector $p_{\mathcal U}$ belongs to it by the Bayes tie, so that complement is $\operatorname{span}(p_{\mathcal U})$.

Now let $v$ be a two-sided feasible direction of $Q_{B_0}$ at $p$. Since $p$ is zero outside $\mathcal U$, two-sided feasibility in the simplex forces $v$ to be zero there. Every active comparison forces $v_{\mathcal U}$ to be orthogonal to \eqref{eq:diffs}. Hence $v_{\mathcal U}=\alpha p_{\mathcal U}$. The simplex constraint gives $0=\1^\top v=\alpha$, because $\1^\top p=1$. Thus $v=0$ and $\mu_{Q_{B_0}}(p)=0$. Since $\|p\|_0=r$, \eqref{eq:ramalower} gives $\CCdim(L^{\rm Jac})\geq r-1$.
\end{proof}

The proof uses positive definiteness only to certify independence of active Bayes comparisons. Positive definiteness by itself would give only an affine upper bound. The factorial witness is what turns the algebra into a lower bound for every convex surrogate.

\section{A two-label falsification certificate}
\label{sec:counterexample}

We now reconcile Theorem~\ref{thm:main} with the purported $s$-score surrogate of \citet{mao2024multilabel}. Their general multi-label logistic loss can be written \[ \widetilde L(h,y) =\sum_{B\subseteq[s]} \Jac(B,y) \log\!\left(\sum_{C\subseteq[s]} e^{\operatorname{score}_h(C)-\operatorname{score}_h(B)}\right), \quad \operatorname{score}_h(B)=\sum_{i=1}^s y_i^B h_i, \tag{20}\label{eq:maoloss} \] where $y_i^B=1$ for $i\in B$ and $-1$ otherwise. The link predicts $\{i:h_i\geq0\}$. Theorem 5.1 and Corollary 5.2 of that work claim a Jaccard excess-risk bound and Bayes consistency.

Take $s=2$, order sets as $\varnothing,\{1\},\{2\},\{1,2\}$, and put \[ p=\left(0,\frac35,\frac1{10},\frac3{10}\right). \tag{21}\label{eq:counterp} \] The expected Jaccard scores of the four reports are \[ a=\E_p[K_{Y,\cdot}] =\left(0,\frac34,\frac14,\frac{13}{20}\right). \tag{22}\label{eq:countera} \] Thus the unique Bayes report is $\{1\}$.

Let \[ s_B(h)=\frac{e^{\operatorname{score}_h(B)}} {\sum_Ce^{\operatorname{score}_h(C)}}. \] The conditional surrogate risk from \eqref{eq:maoloss} is exactly \[ -\sum_B a_B\log s_B(h). \tag{23}\label{eq:crossentropy} \] The probabilities $s_B(h)$ form a product distribution on the two labels. The information projection of $a/\sum_Ba_B$ onto this family matches its label marginals, which here are \[ m_1=\frac{a_{\{1\}}+a_{\{1,2\}}}{\sum_Ba_B}=\frac{28}{33}, \qquad m_2=\frac{a_{\{2\}}+a_{\{1,2\}}}{\sum_Ba_B}=\frac6{11}. \tag{24}\label{eq:marginals} \] Both exceed $1/2$. Consequently the unique finite minimizer is \[ h_1^*=\frac12\log\frac{28}{5}>0, \qquad h_2^*=\frac12\log\frac65>0. \tag{25}\label{eq:hstar} \] The link predicts $\{1,2\}$, whose Jaccard score is $13/20$, instead of the Bayes score $3/4$. Its target regret is $1/10$ at zero surrogate regret. The distribution is supported on nonempty outcomes, so this contradiction is independent of how $\Jac(\varnothing,\varnothing)$ is defined.

The proof gap is equally concrete. The proof of \citet{mao2024multilabel} defines $s^\mu$ by moving probability between two arbitrary label vectors while holding every other coordinate fixed, then asserts that completeness of the $s$ scalar score classes realizes $s^\mu$. But additive softmax probabilities lie on an $s$-dimensional product manifold. With two labels they always satisfy \[ s_\varnothing s_{\{1,2\}}=s_{\{1\}}s_{\{2\}}. \tag{26}\label{eq:productconstraint} \] An arbitrary two-coordinate transfer does not preserve this identity. Coordinatewise completeness makes every $h_i$ attainable, not every point of the $(2^s-1)$-dimensional probability simplex. The same transfer appears in a recent EUM analysis \citep{mohri2026mmo}, but that work addresses population-level generalized metrics rather than the per-instance loss of Theorem~\ref{thm:main}.

\section{Implications and limitations}
\label{sec:discussion}

\paragraph{Why monotone equivalence is insufficient.}
Equation~\eqref{eq:monotone} preserves the ranking of scores for a fixed outcome, but Bayes decisions rank expectations of those scores. Expectation does not commute with a nonlinear monotone transform. The factorial witness shows that this distinction changes convex calibration dimension from quadratic for F1 to exponential for Jaccard. The rank transition in Theorem~\ref{thm:rank} is an algebraic reflection of the same phenomenon.

\paragraph{What the result rules out.}
No loss convex in a subexponential-dimensional report vector, with any link, can be calibrated for expected instance-wise Jaccard on the full conditional simplex. This includes more than coordinate-separable or additive-score losses. The obstruction is distribution-free and link-independent.

\paragraph{What remains possible.}
The theorem does not rule out nonconvex surrogates, restricted conditional families, partial calibration under low noise, or several parallel low-dimensional problems whose total information is exponential. These are genuine escape routes \citep{nueve2024tradeoff}. Nor is calibration dimension itself a running-time lower bound. A method may represent exponentially many scores implicitly and use structured inference. Full-output structured surrogates \citep{mao2023structured} and the generic affine construction match the exponential order.

The target here is the decision-theoretic risk $\E[1-\Jac(h(X),Y)]$. Dataset-level micro IoU, ratios of expected aggregate counts, and empirical utility maximization define different decision problems \citep{nowozin2014iou,mohri2026mmo}. Our theorem does not transfer to those formulations without a separate analysis.

Finally, our upper and lower bounds differ by a factor of two. Determining the exact Jaccard calibration dimension is open. It is also open whether every positive member of \eqref{eq:family} has exponential calibration dimension. The full-rank proof alone cannot establish that claim, and we do not make it.

\bibliography{references}
\bibliographystyle{iclr2027_conference}

\appendix
\section{Calibration-dimension details}
\label{app:ccdim}

This appendix records the exact geometric objects used in
Theorem~\ref{thm:main}. Let $L\in\R_+^{n\times k}$ be a finite loss matrix,
with outcome distributions in $\Delta_n$. For report $t$, write
\[
 Q_t^L=\{q\in\Delta_n:q^\top L_{\cdot,t}
            \leq q^\top L_{\cdot,t'}\text{ for every }t'\in[k]\}.
 \tag{27}\label{eq:triggerapp}
\]
For a convex set $Q\subseteq\R^n$ and $p\in Q$, its cone of feasible
directions is
\[
 F_Q(p)=\{v:\text{there is }\epsilon_0>0\text{ such that }
 p+\epsilon v\in Q\text{ for every }\epsilon\in(0,\epsilon_0)\}.
\]
The feasible subspace is $F_Q(p)\cap(-F_Q(p))$, and its dimension is
$\mu_Q(p)$. Theorem 16 of \citet{ramaswamy2016ccdim} gives
\[
 \CCdim(L)\geq\|p\|_0-\mu_{Q_t^L}(p)-1
 \quad\text{for every }p\in Q_t^L.
 \tag{28}\label{eq:ramatheorem}
\]
It applies to the calibration definition used in their paper and does not
require the conditional surrogate risk to attain its infimum.

For completeness, their affine upper bound states
\[
 \CCdim(L)\leq\affdim(L),
 \tag{29}\label{eq:affupper}
\]
where $\affdim(L)$ is the affine dimension of the set of loss columns. One
way to view the construction is to choose affine coordinates for the loss
columns, estimate the corresponding conditional expectations with a convex
quadratic loss, and decode by minimizing the reconstructed target risk.

\section{Proof details for the rank transition}
\label{app:rank}

We prove the positive-$\lambda$ part of Theorem~\ref{thm:rank} with all
interchanges explicit. Let $\mathcal Y_+=2^{[s]}\setminus\{\varnothing\}$.
For $A,B\in\mathcal Y_+$, abbreviate
\[
 x=|A\cap B|,\qquad a=|A|,\qquad b=|B|.
\]
Because $0\leq\lambda<2$ and $x\leq\min(a,b)$, the denominator in
\eqref{eq:family} is positive. The elementary identity
\[
 \frac{x}{a+b-\lambda x}
 =\int_0^\infty xe^{-t(a+b-\lambda x)}dt
 \tag{30}\label{eq:laplaceapp}
\]
gives \eqref{eq:integral}.

For incidence vectors $\chi_A,\chi_B\in\{0,1\}^s$,
\[
 x^r=\langle\chi_A^{\otimes r},\chi_B^{\otimes r}\rangle.
 \tag{31}\label{eq:tensoridentity}
\]
Moreover,
\[
 xe^{\lambda tx}
 =\sum_{r=1}^\infty
 \frac{(\lambda t)^{r-1}}{(r-1)!}x^r.
 \tag{32}\label{eq:seriesapp}
\]
All coefficients in this series are nonnegative. Since the set family is
finite, Tonelli's theorem applies after separating the positive and negative
parts of the finite coefficient vector, or equivalently after first expanding
the finite quadratic form. Equations~\eqref{eq:laplaceapp}--\eqref{eq:seriesapp}
yield
\begin{align}
 \sum_{A,B\in\mathcal Y_+}c_Ac_BS_\lambda(A,B)
 &=(2-\lambda)\int_0^\infty
 \sum_{r=1}^\infty\frac{(\lambda t)^{r-1}}{(r-1)!}
 \left\langle
 \sum_Ac_Ae^{-t|A|}\chi_A^{\otimes r},
 \sum_Bc_Be^{-t|B|}\chi_B^{\otimes r}
 \right\rangle dt\\
 &=(2-\lambda)\int_0^\infty
 \sum_{r=1}^\infty\frac{(\lambda t)^{r-1}}{(r-1)!}
 \left\|\sum_Ac_Ae^{-t|A|}\chi_A^{\otimes r}\right\|_2^2dt.
 \tag{33}\label{eq:quadraticapp}
\end{align}
The integrand is continuous and nonnegative. If the integral is zero, the
integrand is zero for every $t>0$. Since $\lambda>0$, every tensor norm in
the sum is zero.

Fix a nonempty $T=\{i_1,\ldots,i_r\}\subseteq[s]$. The
$(i_1,\ldots,i_r)$ coordinate of the $r$th tensor sum is
\[
 \sum_{A\in\mathcal Y_+}c_Ae^{-t|A|}
 \prod_{j=1}^r\chi_A(i_j)
 =\sum_{A\supseteq T}c_Ae^{-t|A|}.
 \tag{34}\label{eq:zetacoord}
\]
Fix any $t>0$ and put $d_A=c_Ae^{-t|A|}$. Equation~\eqref{eq:zetacoord}
says that the upper-zeta transform $T\mapsto\sum_{A\supseteq T}d_A$ is zero
on the poset of nonempty subsets. Ordered by decreasing cardinality, its
matrix is triangular with ones on the diagonal. Hence $d_A=0$ for every
$A$, and therefore $c_A=0$. The nonempty block is strictly positive
definite. The empty set has score zero against every nonempty set and score
one against itself, so it contributes a separate positive $1\times1$ block.

At $\lambda=0$, only the $r=1$ term remains in
\eqref{eq:seriesapp}. The variation of $e^{-t|A|}$ with cardinality allows
more than $s$ rank, but not full rank. \citet{zhang2026exact} prove the exact
value $s^2-s+2$, including the empty-set block.

Finally, if a square matrix $K$ of size $N$ is nonsingular, its $N$ columns
are linearly independent and thus affinely independent. The affine map
$k_B\mapsto\1-k_B$ preserves affine independence, so the loss columns of
$\1\1^\top-K$ have affine dimension $N-1$. This proves
\eqref{eq:affdim}.

\section{The factorial tie identity}
\label{app:tie}

We give a full proof of Lemma~\ref{lem:tie}. Recall that $O$ contains the
$n=s-1$ optional labels. For $C\subseteq O$, define $W(C)$ by
\eqref{eq:W}. We first show that $W(C)$ is invariant under adding an element.

Fix $j\in O\setminus C$. Pair the outcome index $D=E$ with
$D=E\cup\{j\}$ for each $E\subseteq O\setminus\{j\}$. Put
\[
 a=|E\cap C|,\qquad b=|E\setminus C|,\qquad c=|C|.
\]
For $D=E$, adding $j$ to the report leaves the intersection unchanged and
increases the union size by one. Its contribution to
$W(C\cup\{j\})-W(C)$ is
\[
 -\frac{a+1}{(a+b)!(c+b+1)(c+b+2)}.
 \tag{35}\label{eq:nojinD}
\]
For $D=E\cup\{j\}$, adding $j$ to the report leaves the union unchanged and
increases the intersection by one. Its contribution is
\[
 \frac1{(a+b+1)!(c+b+2)}.
 \tag{36}\label{eq:jinD}
\]
Their sum is \eqref{eq:increment}.

For a fixed $b$, the number of choices of $E$ having a given $a$ is
$\binom ca\binom{n-c-1}{b}$. The second binomial coefficient and the factor
$1/(c+b+2)$ do not depend on $a$. It remains to prove
\eqref{eq:identity}. Multiply its left-hand side by $c+b+1$ and use
$c+b+1=(c-a)+(a+b+1)$:
\begin{align}
 (c+b+1)\sum_{a=0}^c\binom ca\frac1{(a+b+1)!}
 &=\sum_{a=0}^c\binom ca\frac1{(a+b)!}
   +\sum_{a=0}^c(c-a)\binom ca\frac1{(a+b+1)!}\\
 &=\sum_{a=0}^c\binom ca\frac1{(a+b)!}
   +c\sum_{a=0}^{c-1}\binom{c-1}{a}\frac1{(a+b+1)!}\\
 &=\sum_{a=0}^c\binom ca\frac1{(a+b)!}
   +\sum_{a=0}^ca\binom ca\frac1{(a+b)!}\\
 &=\sum_{a=0}^c(a+1)\binom ca\frac1{(a+b)!}.
 \tag{37}\label{eq:identityproof}
\end{align}
In the third line we shifted the index and used
$a\binom ca=c\binom{c-1}{a-1}$. Hence every fixed-$b$ sum of paired
increments is zero, and $W(C\cup\{j\})=W(C)$.

Repeatedly removing elements from $C$ gives $W(C)=W(\varnothing)$. Directly,
\[
 W(\varnothing)
 =\sum_{D\subseteq O}\frac1{|D|!(1+|D|)}
 =\sum_{d=0}^n\binom nd\frac1{(d+1)!},
 \tag{38}\label{eq:wempty}
\]
which proves \eqref{eq:constant}. After division by $Z_n$, all reports in
$\mathcal U$ therefore tie.

Now take any report $B$ with $\star\notin B$ and put
$B^+=B\cup\{\star\}$. Every outcome $A$ in the support of $p$ contains
$\star$. Therefore
\[
 |B^+\cup A|=|B\cup A|,
 \qquad
 |B^+\cap A|=|B\cap A|+1.
\]
The Jaccard score of $B^+$ is strictly larger for every such $A$. Since all
support probabilities are positive, $B$ is strictly suboptimal. Thus exactly
$\mathcal U$ is Bayes optimal.

\section{Feasible-subspace calculation}
\label{app:feasible}

We expand the geometric step in Theorem~\ref{thm:main}. Let
$K_{\mathcal U,\mathcal U}$ be the $r\times r$ principal Jaccard block, where
$r=2^{s-1}$. It is positive definite by Theorem~\ref{thm:rank}. Fix
$B_0\in\mathcal U$ and form the matrix whose $r-1$ columns are
\[
 d_B=K_{\mathcal U,B}-K_{\mathcal U,B_0},
 \qquad B\in\mathcal U\setminus\{B_0\}.
 \tag{39}\label{eq:dBapp}
\]
These columns are independent. Indeed, if
$\sum_{B\ne B_0}\alpha_Bd_B=0$, define coefficients $\beta_B=\alpha_B$ for
$B\ne B_0$ and $\beta_{B_0}=-\sum_{B\ne B_0}\alpha_B$. Then
$K_{\mathcal U,\mathcal U}\beta=0$. Nonsingularity gives $\beta=0$, hence
every $\alpha_B=0$. Thus the matrix $D$ of differences has rank $r-1$.

The tie identity says $D^\top p_{\mathcal U}=0$. Since
$\dim\ker(D^\top)=1$ and $p_{\mathcal U}\ne0$,
\[
 \ker(D^\top)=\operatorname{span}(p_{\mathcal U}).
 \tag{40}\label{eq:kernelapp}
\]

Let $v\in F_{Q_{B_0}}(p)\cap(-F_{Q_{B_0}}(p))$. If $A\notin\mathcal U$,
then $p_A=0$. Feasibility of both $v$ and $-v$ in the nonnegative orthant
forces $v_A=0$. For every $B\in\mathcal U$, the target risks of $B$ and
$B_0$ tie at $p$. The corresponding inequality in $Q_{B_0}$ is active.
Feasibility in both directions forces its directional derivative to be zero:
\[
 v_{\mathcal U}^\top
 (L^{\rm Jac}_{\mathcal U,B}-L^{\rm Jac}_{\mathcal U,B_0})=0.
 \tag{41}\label{eq:activeapp}
\]
Since loss differences are negative score differences,
\eqref{eq:activeapp} is $D^\top v_{\mathcal U}=0$. By
\eqref{eq:kernelapp}, $v_{\mathcal U}=\alpha p_{\mathcal U}$. Finally,
two-sided feasibility in the affine hull of the simplex requires
$\1^\top v=0$. Because $\1^\top p=1$, this implies $\alpha=0$. The feasible
subspace is $\{0\}$ and has dimension zero.

Notice that constraints against reports outside $\mathcal U$ were not needed.
They are strict at $p$ by Lemma~\ref{lem:tie} and therefore impose no
additional equality on a two-sided feasible direction.

\section{Exact analysis of the logistic counterexample}
\label{app:counterexample}

For the set order
$\varnothing,\{1\},\{2\},\{1,2\}$, the Jaccard score matrix is
\[
 K=
 \begin{pmatrix}
 1&0&0&0\\
 0&1&0&1/2\\
 0&0&1&1/2\\
 0&1/2&1/2&1
 \end{pmatrix}.
 \tag{42}\label{eq:Ktwo}
\]
Multiplying \eqref{eq:counterp} by this matrix gives
\eqref{eq:countera}. The unique maximum $3/4$ occurs at report $\{1\}$.

We next derive the conditional surrogate risk. Let
$r_B=\operatorname{score}_h(B)$. Each term in \eqref{eq:maoloss} satisfies
\[
 \log\sum_Ce^{r_C-r_B}
 =\log\sum_Ce^{r_C}-r_B=-\log s_B(h).
 \tag{43}\label{eq:logsoftmaxapp}
\]
Taking expectation over $Y$ and interchanging two finite sums gives
\[
 \E_p\widetilde L(h,Y)
 =-\sum_B\E_p[\Jac(B,Y)]\log s_B(h)
 =-\sum_Ba_B\log s_B(h).
 \tag{44}\label{eq:riskapp}
\]

Because $r_B=\sum_i y_i^Bh_i$, $s_B(h)$ is the product of two Bernoulli
probabilities. If $q_B=a_B/A$ with $A=\sum_Ba_B=33/20$, then
\eqref{eq:riskapp} is $A$ times the cross entropy from $q$ to this product
family. Writing $\pi_i=\Pr_{s(h)}(i\in B)$, the part depending on $h$ is
\[
 A\sum_{i=1}^2[-m_i\log\pi_i-(1-m_i)\log(1-\pi_i)],
 \tag{45}\label{eq:binaryce}
\]
up to a constant independent of $h$. Each binary cross entropy is strictly
convex in its logit and uniquely minimized at $\pi_i=m_i$. Direct arithmetic
gives
\[
 A=\frac{33}{20},\qquad
 m_1=\frac{3/4+13/20}{33/20}=\frac{28}{33},\qquad
 m_2=\frac{1/4+13/20}{33/20}=\frac6{11}.
 \tag{46}\label{eq:marginalarithmetic}
\]
Under the $\{-1,1\}$ score encoding,
$\pi_i=e^{h_i}/(e^{h_i}+e^{-h_i})$. Solving for $h_i$ yields
\eqref{eq:hstar}. Both coordinates are positive, so the sign link returns
$\{1,2\}$. Its target loss exceeds the Bayes loss by
\[
 \left(1-\frac{13}{20}\right)-\left(1-\frac34\right)=\frac1{10}.
 \tag{47}\label{eq:regretapp}
\]
The conditional surrogate risk is minimized exactly, so its regret is zero.

For reference, the product constraint \eqref{eq:productconstraint} follows
immediately from the four unnormalized weights
\[
 e^{-h_1-h_2},\quad e^{h_1-h_2},\quad
 e^{-h_1+h_2},\quad e^{h_1+h_2}.
\]
Both cross-products equal one before applying the common softmax normalizer.

\section{Finite exact checks}
\label{app:checks}

The proof is symbolic, but small cases provide useful falsification checks.
For $s=2,3,4,5$, exact rational calculations for the witness
\eqref{eq:witness} give
\[
\begin{array}{c|cccc}
s & |\mathcal U| & \text{tied score} & \text{best score outside }\mathcal U
  & \operatorname{rank}(K_{\mathcal U,\mathcal U})\\ \hline
2&2&3/4&1/4&2\\
3&4&13/21&2/7&4\\
4&8&73/136&39/136&8\\
5&16&501/1045&292/1045&16
\end{array}
\tag{48}\label{eq:smallchecks}
\]

The calculation enumerates all $2^s$ outcomes and reports using rational
arithmetic, checks equality of the scores over $\mathcal U$, strict
separation outside $\mathcal U$, and exact rank of the rational principal
matrix. The two-label certificate separately checks
\[
 a=(0,3/4,1/4,13/20),\quad
 (m_1,m_2)=(28/33,6/11),\quad
 \text{regret}=1/10.
\]
No numerical approximation is used in these checks.


\end{document}
